\chapter{Introduction} \label{ch:introduction}

Bitcoin's custody model places the entire burden of key security on the asset holder.
A compromised private key results in irreversible fund loss, with no recourse mechanism
built into the protocol. Covenant-based vaults address this by interposing a time-delayed
spending path between key compromise and fund movement: when the hot key initiates a
withdrawal, a configurable delay window allows a watchtower to detect the unauthorized
spend and trigger emergency recovery to a cold storage address~\cite{SHMB20, MES16}.

Four concrete covenant proposals have emerged as candidates for Bitcoin activation:
\opctv{}~\cite{BIP119}, \opccv{}~\cite{BIP443}, \opvault{}~\cite{BIP345}, and
\catcsfs{}~\cite{BIP347, BIP348}. Each enforces the vault lifecycle through a different
mechanism. CTV pre-commits to the exact transaction structure at vault creation time.
CCV validates contract state transitions dynamically at spend time. OP\_VAULT provides
dedicated vault opcodes with built-in recovery authorization. CAT+CSFS composes two
general-purpose opcodes to achieve transaction introspection via dual signature
verification. A fifth design, built with the Simplicity language on the Elements
sidechain~\cite{OC17}, demonstrates what covenant enforcement looks like with full
transaction introspection and no soft-fork dependency.

Despite active community deliberation on which proposal to activate---including a
66-developer open letter requesting Bitcoin Core prioritize CTV+CSFS
review~\cite{CTVCSFS25}, calls for comprehensive analysis before
activation~\cite{Poinsot25, Towns25}, and the subsequent withdrawal of BIP-345
in favor of BIP-443~\cite{OB25}---no published work compares these designs
empirically. The BIP
specifications describe each opcode in isolation. Qualitative comparison matrices
exist~\cite{covenants_info}, but they contain no transaction sizes, no attack cost
functions, and no fee-environment analysis. The only quantitative estimate in the
literature is Harding's back-of-envelope calculation of watchtower splitting capacity
at approximately 3{,}000 splits per block~\cite{Har24}, based on assumed---not
measured---transaction sizes.

This thesis fills that gap. We build all four Bitcoin vault implementations, instrument
them through a uniform adapter interface, and run 15 experiments on regtest measuring
transaction costs, security properties, and capability differences. We additionally
build two original vault implementations---a CAT+CSFS vault using the explicit
\texttt{OP\_CHECKSIGFROMSTACK} (BIP-348) dual verification pattern, and a Simplicity
vault on Elements using jet-based covenant enforcement---and develop Alloy bounded
model checking specifications for all five designs.

\section{Main Results} \label{se:main_results}

Two results constrain the design space of covenant vaults.

\paragraph{Griefing--safety impossibility.}
Recovery in a covenant vault is either permissionless (anyone can invoke it, as in CCV)
or key-gated (requiring a specific key, as in OP\_VAULT). Permissionless recovery
maximizes fund safety under key loss---the owner never loses the ability to recover---but
enables denial-of-service griefing by any third party who can broadcast a recovery
transaction. Key-gated recovery eliminates griefing but introduces key-loss risk: if the
recovery authorization key is lost, funds become permanently unrecoverable. We formalize
this as Proposition~\ref{thm:griefing_safety} in Chapter~\ref{ch:security}: no covenant vault
simultaneously achieves permissionless recovery and griefing resistance. This tradeoff is
a design constraint, not an implementation artifact.

\paragraph{Per-threat worst-case characterization.}
CTV vaults are vulnerable to descendant-chain fee pinning, where an attacker constructs
25 low-value transactions off the unvault output to block the watchtower's
child-pays-for-parent recovery. CCV and OP\_VAULT vaults are vulnerable to watchtower
exhaustion, where an attacker repeatedly splits the vault into smaller UTXOs until
individual recovery transactions become uneconomical. Rather than aggregating
disparate attack outcomes (theft versus griefing) into a single ``security ranking'',
we formalise a per-strategy defender-loss metric~$\mathcal{L}$ and characterise, threat
by threat, how~$\mathcal{L}$ varies across designs with the fee rate~$r$ and the
defender's batching strategy. The per-threat analysis surfaces a combinatorial
impossibility---no design minimises $\mathcal{L}$ across all threats and fee
rates---that is a direct corollary of the orderings, not a standalone claim about
``safest'' designs.

\section{Contributions} \label{se:contributions}

This work makes seven contributions:

\begin{enumerate}
\item \textbf{First empirical comparison of the four Bitcoin Script extensions.}
  Regtest-measured transaction sizes for CTV, CCV, OP\_VAULT, and CAT+CSFS under a
  uniform adapter interface, spanning 15~experiments and 11~threat models. We correct
  prior hand-estimates of OP\_VAULT trigger transaction size from approximately
  200\vB{} to 292\vB{} (a 46\% error), revealing that OP\_VAULT's lifecycle is 36\% more
  expensive than CCV's.

\item \textbf{Formal adversary-advantage framework.}
  A defender-loss metric~$\mathcal{L}(\mathrm{TM}, D, V, r, b)$ and adversary-advantage
  metric~$\mathrm{Adv}$ (Definitions~\ref{def:strategy_costs}
  and~\ref{def:loss_adv}) that uniformly compare theft, griefing, and dust-lockup
  outcomes across threat models, fee rates, and defender batching strategies.

\item \textbf{Per-threat worst-case characterisation.}
  For each of TM1--TM11, a closed-form or bounded expression for~$\mathcal{L}$
  across designs as a function of fee rate, including the dust-threshold rate
  $r^\dagger$ for per-UTXO recovery-cost scaling (Section~\ref{se:thm_inversion}).

\item \textbf{Griefing--safety impossibility (Proposition~\ref{thm:griefing_safety}).}
  No covenant vault achieves both permissionless recovery and griefing resistance.

\item \textbf{Structural design-space decomposition
  (Proposition~\ref{thm:fee_inversion}).}
  A four-dimensional decomposition of covenant-vault design choices
  ($f$: fee model; $a$: amount flexibility; $g$: recovery gating;
  $b$: recovery binding). Each dimension has a deterministic
  feature-to-vulnerability implication
  (Lemmas~\ref{lem:anchor_pinning}--\ref{lem:unbound_theft}). No
  proposed Bitcoin Script extension occupies the dominant corner
  $(f_{\neq\text{anchor}}, a_{\text{atomic}}, g_{\text{key}},
  b_{\text{bound}})$; each falls short by at least one dimension with
  a corresponding vulnerability surface.

\item \textbf{Batched-defender measurements and ordering-flip finding.}
  On-chain measurements of BIP-345 and BIP-443 batched recovery yield
  $(\alpha, \beta) = (54, 68)$\vB{} for CCV and $(154, 92)$\vB{} for OP\_VAULT. The
  batched dust-threshold rate $r^\dagger$ rises by $1.8\times$ (CCV) and $2.7\times$
  (OP\_VAULT), and the CCV/OP\_VAULT ordering under per-UTXO recovery-cost scaling reverses
  (Section~\ref{ss:batched_defender}).

\item \textbf{Two original vault implementations and Alloy formal verification.}
  A CAT+CSFS vault using the explicit \texttt{OP\_CHECKSIGFROMSTACK} (BIP-348) opcode
  for dual signature verification, distinct from prior OP\_CAT-only vaults that
  achieve introspection via the Schnorr G-trick without CSFS~\cite{Poe21,
  Rij24}. Additionally, a Simplicity vault using \texttt{jet::outputs\_hash()} on
  Elements regtest via the Simplex SDK, treated throughout as a cross-substrate
  reference point rather than a peer of the four Bitcoin proposals. Five Alloy
  models verify fund conservation, recovery liveness, and no-extraction properties.
\end{enumerate}

\section{Scope and Limitations} \label{se:scope}

All experiments run on Bitcoin Core regtest (or Elements regtest for Simplicity).
Transaction virtual size (vsize) is the primary metric; it is structurally deterministic
and identical on regtest and mainnet. Fee amounts are computed as projections
($\text{fee} = \text{vsize} \times r$, where $r$ is an exogenous fee rate) rather than
measured directly, because regtest has no fee market. Temporal dynamics of recovery
races---where confirmation latency determines whether the watchtower or the attacker
wins---are not captured, because regtest mines blocks on demand.

The four Bitcoin Script extension proposals are tested using their reference
implementations: \texttt{simple-ctv-vault}~\cite{BIP119}, \texttt{pymatt}~\cite{BIP443},
\texttt{opvault-demo}~\cite{BIP345}, and our \texttt{cat-csfs-vault}. Where a reference
implementation has a bug or missing feature, we distinguish this from a design-level
property of the underlying BIP. Simplicity on Elements is included throughout as a
cross-substrate reference point rather than a peer of the four Bitcoin proposals.
Elements is a federated sidechain with functionary block production rather than
proof-of-work, and the threat taxonomy does not transfer identically across that
substrate boundary. Our formal results
(Proposition~\ref{thm:fee_inversion}, Proposition~\ref{thm:griefing_safety}) therefore
range over the four Bitcoin Script extensions only; Simplicity measurements are
presented separately.

Protocol-layer constructs such as Ark and BitVM are excluded. Both operate above the
opcode level---Ark uses an operator-mediated virtual UTXO model, BitVM uses off-chain
fraud proofs---and require modeling interactive communication rounds that fall outside
our single-transaction covenant framework.

\section{Thesis Organization} \label{se:organization}

Chapter~\ref{ch:background} surveys the four Bitcoin Script extension proposals and
introduces Simplicity as an alternative script primitive on Elements, defining the
vault lifecycle model and threat taxonomy used throughout.
Chapter~\ref{ch:methodology} describes the measurement framework, adapter architecture,
and regtest validity analysis. Chapter~\ref{ch:results} presents the empirical
measurements: lifecycle costs, batching efficiency, revault amplification, and the
corrections to prior estimates. Chapter~\ref{ch:security} contains the core security
analysis, including Propositions~\ref{thm:griefing_safety}
and~\ref{thm:fee_inversion},
parameterized attack cost functions, and the full threat model matrix.
Chapter~\ref{ch:verification} describes the Alloy formal verification methodology and
results. Chapter~\ref{ch:discussion} discusses limitations, deployment guidance, and
open problems. Chapter~\ref{ch:related} positions this work relative to prior vault
security research. Chapter~\ref{ch:conclusion} summarizes findings and future directions.

\begin{table}[t]
\centering
\caption{Five covenant vault designs compared in this work. ``BIP status'' refers
to the proposal's standing in the Bitcoin Improvement Proposal process, not
deployment on mainnet---none of the Bitcoin proposals are activated on mainnet
as of April 2026.}
\label{tab:designs_overview}
\begin{tabular}{@{}llllc@{}}
\toprule
Design & BIP & Enforcement mechanism & Chain & BIP status \\
\midrule
CTV        & 119       & Template hash commitment   & Bitcoin  & Proposed \\
CCV        & 443       & Contract state verification & Bitcoin  & Proposed \\
OP\_VAULT  & 345       & Dedicated vault opcodes     & Bitcoin  & Withdrawn \\
CAT+CSFS   & 347 + 348 & Dual signature introspection & Bitcoin & Proposed \\
Simplicity & ---       & Jet-based introspection     & Elements & Deployed \\
\bottomrule
\end{tabular}
\end{table}
